% Generated draft: requires mathematical and source review.
\documentclass[11pt]{article}
\usepackage[margin=25mm]{geometry}
\usepackage{fontspec}
\setmainfont{Latin Modern Roman}
\usepackage{amsmath,amssymb,mathtools,unicode-math}
\setmathfont{Latin Modern Math}
\usepackage{enumitem,needspace,xurl,hyperref}
\setlength{\parindent}{0pt}
\setlength{\parskip}{5pt}
\setlength{\emergencystretch}{3em}
\newcommand{\N}{\mathbb{N}}
\newcommand{\Z}{\mathbb{Z}}
\newcommand{\Q}{\mathbb{Q}}
\newcommand{\R}{\mathbb{R}}
\newcommand{\C}{\mathbb{C}}
\newcommand{\F}{\mathbb{F}}
\newcommand{\A}{\mathbb{A}}
\newcommand{\PP}{\mathbb{P}}
\newcommand{\E}{\mathbb{E}}
\newcommand{\eps}{\varepsilon}
\newcommand{\dd}{\,\mathrm{d}}
\begin{document}
\section*{UnsolvedMath ID: 800734}
\textbf{Source problem code:} catalogue entry
\subsection{Definitions and mathematical statement}
A topological group is a group \(G\) equipped with a topology for which multiplication \(G\times G\to G\) and inversion \(G\to G\) are continuous. The supplied record does not specify a separation axiom; the intended convention requires review. The group is non-discrete if its topology is not discrete. A subset \(A\subseteq G\) is nowhere dense if
\[
\operatorname{int}_G\!\left(\overline{A}^{\,G}\right)=\varnothing.
\]
Its algebraically generated subgroup is
\[
\langle A\rangle
=\{a_1^{\varepsilon_1}\cdots a_n^{\varepsilon_n}: n\in\Z_{\geq 0},\ a_i\in A,\ \varepsilon_i\in\{-1,1\}\},
\]
where the empty product is the identity.

The question supplied as Problem 8, attributed to Protasov, is: does every non-discrete topological group \(G\) admit a subset \(A\subseteq G\) such that
\[
\operatorname{int}_G\!\left(\overline{A}^{\,G}\right)=\varnothing
\qquad\text{and}\qquad
\langle A\rangle=G?
\]

The record also includes a parenthetical assertion explicitly referring to \emph{Problem 6}: its answer is reported as negative, with the following conditional example attributed to Proposition 13.8 of the Protasov--Banakh monograph cited below. Under the continuum hypothesis \(2^{\aleph_0}=\aleph_1\), the product topological group \((\R^{\omega_1},+)\), where \(\omega_1\) is the first uncountable ordinal, is reported to contain a dense non-separable subgroup \(H\) such that every meager subset \(M\subseteq H\) generates a subgroup \(\langle M\rangle\) that is metrizable and separable in the subspace topology. Here a meager subset is a countable union of nowhere dense subsets. This is a supplied source assertion, not an independently verified theorem in this draft. The reference to Problem 6, its relationship to the displayed Problem 8 question, and the registry's supplied open status remain unresolved.
\subsection{Short English statement}
The question asks whether every non-discrete topological group can be generated, using finite products and inverses, by a nowhere dense subset. The supplied record also reports a conditional negative example under the continuum hypothesis, but labels that discussion as concerning Problem 6. This numbering and status issue requires review.
\subsection{Sources}
\begin{itemize}
\item Local UnsolvedMath problem registry, permanent ID \(800734\), exact source problem code \texttt{catalogue entry}. The supplied record labels the problem \texttt{open}; this is a registry label, not an independently verified assessment.
\item Banakh and Choban, \emph{Some Open Problems in Topological Algebra} (2012), Problem 8. Supplied source URL: \url{https://arxiv.org/abs/1202.1619}. The record attributes the problem to Protasov. Its provenance reports access on 2026-10-16 and extraction by \texttt{arxiv\_latex}; these are supplied provenance details, not a claim of source access for this draft.
\item I. Protasov and T. Banakh, \emph{Ball stuctures and colorings of graphs and groups}, Mat. Stud. Monograph Series, volume 11, VNTL Publ., 2003, 148 pages, Proposition 13.8. This citation and title spelling are reproduced from the supplied parenthetical concerning Problem 6; the proposition has not been independently checked.
\end{itemize}
\end{document}
